\documentclass[10pt]{article}

\usepackage[utf8]{inputenc}

\usepackage[T1]{fontenc}

\usepackage{amsmath}

\usepackage{amsfonts}

\usepackage{amssymb}

\usepackage[version=4]{mhchem}

\usepackage{stmaryrd}

\usepackage{bbold}

\usepackage{hyperref}

\hypersetup{colorlinks=true, linkcolor=blue, filecolor=magenta, urlcolor=cyan,}

\urlstyle{same}



\begin{document}

С помощью С. о. описываются многие краевые задачи математич. физики.



Лum.: [1] Л ю с т е р н ик Л. А., С о б о л е в В. И., Элементы функционального анализа, 2 изд., М., 1965; [2] А х ие з е р Н. И., Г л а з м а н И. М., Теория линейных операторов в гильбертовом пространстве, М., $1966 ;$ [3] Р и с с Ф., анализу, пер. с франц., 2 изд., М., $1979 . \quad$ В. И. Соболев.



САРДА ТЕОРЕМА: пусть $f: M \rightarrow N$ - отображение класса $C^{r}$ многообразий $M$ и $N$ размерности $m$ и $n$ соответственно; если $r>\max (0, m-n)$, то критические значения $f$ образуют множество меры нуль. Множество же регулярных значений оказывается массивным и всюду плотным. Установлена А. Сардом [1].



JIum.: [1] S a r d A., "Bull. Amer. Math. Soc.», 1942, v. 48, р. С83-90 ННСИРОВАННОЕ КОЛБЦО Л е в о е (п р ав о е) - колыц, над к-рым все левые (правые) модули сбалансированы. Кольцо сбалансировано слева тогда и только тогда, когда все его факторкольца суть $Q F-1-$ к о л ц ц, т. е. все точные левые модули над ними сбалансированы. В частности, кольцо сбалансировано, если все эти факторкольца квазифробениусовы. Всякое С. к. разлагается в прямую сумму однорядного кольца и колец матриц над локальными кольцами специального типа. Јюбое С. к. полусовершенно. Нётерово С. к. оказывается артиновым.



Лит.: [1] Итоги науки и техники. Алгебра. Топология. Геометрия, т. 19, М., 1981, с. 31-134; [2] Ф е й с К., Алгебра: кольца, модули и категории, пер. с англ., т. 1-2, М., $1977-79$.



СБАЛАНСИРОВАННЫИ моду ль - модуль $M$ такой, что естественный кольцевой гомоморфизм $\varphi: R \rightarrow$ End $_{\text {End }_{R}} M$, в случае правого модуля определяемый равенством $\varphi(r)(m)=m r$ для любых $r \in R$ и $m \in M$, сюръективен. Модуль $P$ над кольцом $R$ оказывается образующим категории $R$-модулей тогда и только тогда, когда $P$ есть С. м. как $R$-модуль, проективен и конечно порожден как End $R_{R} P$-модуль. Лит.: [1] Ф е й с К., Алгебра: кольца, модули и категории,

пер. с англ., т. $1-2$, М., $1977-79$. СВЕРТКА ф у н ц ц и й $f(x)$ и $g(x)$, п р и н а д л е ж а и х $L(-\infty, \infty)$, - фунция $h(x)$, определяемая равенством



\[

h(x)=\int_{-\infty}^{\infty} f(x-y) g(y) d y=\int_{-\infty}^{\infty} f(y) g(x-y) d y

\]



П обозначаемая символом $(f * g)(x)$. Функция $f * g$ определена почти всюду и также принадлежит $L(-\infty, \infty)$. Свертка обладает основными свойствами операции умножения, а именно: $f * g=g * f$,



\[

\begin{aligned}

\left(\alpha_{1} f_{1}+\alpha_{2} f_{2}\right) * g= & \alpha_{1}\left(f_{1} * g\right)+\alpha_{2}\left(f_{2} * g\right), \alpha_{1}, \alpha_{2} \in \mathbb{R}, \\

& (f * g) * h=f *(g * h)

\end{aligned}

\]



для любых трех функций из $L(-\infty, \infty)$. Поэтому $L(-\infty, \infty)$ с обычным слож́нением и умножением на число, с операцией С. в качестве умножения әлементов из $L(-\infty, \infty)$ и с нормой



\[

\|f\|=\int_{-\infty}^{\infty}|f(x)| d x

\]



превращается в банахову алаебру (при такой норме $\left.\left\|\mathrm{f}^{*} g\right\| \leqslant\|f\| \cdot\|g\|\right)$. Если $F[f]$ - преобразование Фурье функции $f$, то



\[

F[f * g]=\sqrt{2 \pi} F[f] F(g)

\]



и это используется Іри решении ряда шрикладных задач.



Так, если задача сведена к интегральному уравнению вида



\[

f(x)=g(x)+\int_{-\infty}^{\infty} K(x-y) f(y) d y

\]



где



\[

g(x) \in L_{2}(-\infty, \infty), K(x) \in L(-\infty, \infty) \text {, }

\][



\textbackslash sup \_\{x\}|F\href{x}{K}| \textbackslash leqslant \textbackslash frac\{1\}\{\textbackslash sqrt\{2 \textbackslash pi\}\}

]



то в предположении, что $f(x) \subset L(-\infty, \infty)$, применяя к уравнению (*) преобразование Фурье, получают



\[

F[f]=F[g]+\sqrt{2 \pi} F[f] F[K],

\]



откуда



\[

F[f]=\frac{F[g]}{1-\sqrt{2 \pi} F[K]}

\]



и обратное преобразование Фурье приводит к решению



\[

f(x)=\frac{1}{\sqrt{2 \pi}} \int_{-\infty}^{\infty} \frac{F[g](\zeta) e^{-i \zeta x}}{1-\sqrt{2 \pi} F[K](\zeta)} d \zeta

\]



уравнения (*).



Свойства С. функций находят важные приложения в теории вероятностей. Если $f(x)$ и $g(x)$ являются плотностями вероятности независимых случайных величин $X$ и $Y$, то С. $(f * g) x$ есть плотность вероятности случайной величины $X+Y$.



Операция C. распространяется на обобщенные бункцuu. Если $f$ и $g$ - обобщенные функции, из к-рых по крайней мере одна имеет компактный носитель, и $\varphi(x)$ принадлежит пространству основных функций, то $f * g$ определяется равенством



\[

\langle f * g, \varphi\rangle=\langle f(x) \times g(y), \varphi(x+y)\rangle,

\]



где $f(x) \times g(y)$ - прямое произведение обобщенных функций $f$ и $g$, т. е. функционал в пространстве основных функций двух независимых переменных такой, что



\[

\langle f(x) \times g(x), u(x, y)\rangle=\langle f(x),\langle g(y), u(x, y)\rangle\rangle

\]



для любой финитной бесконечно диффференцируемой функции $u(x, y)$.



C. обобщенных функций также обладает свойством коммутативности, линейности по каждому аргументу, а если по крайней мере две из трех обобщенных функций имеют компактные носители, то и свойством ассоциативности. Справедливы равенства:



\[

D^{\alpha}(f * g)=D^{\alpha} f * g=f * D^{\alpha} g,

\]



где $D$ оператор дифференцирования и $\alpha-$ любой мультииндекс, $\left(D^{\alpha} \delta\right) * f=D^{\alpha} f$, в частности $\delta * f=f$, где $\delta$ - дельта-функция, и если $f_{n}, n=1,2, \ldots$, - обобщенные функции такие, что $f_{n} \rightarrow f_{0}$ и существует компакт



\[

\begin{gathered}

K \supset \operatorname{supp} f_{n}, n=1,2, \ldots, \\

f_{n} * g \rightarrow f_{0} * g .

\end{gathered}

\]



Наконец, если $g$ - финитная обобщенная функция и $f$ - обобщенная функция медленного роста, то к $f * g$ применимо преобразование Фурье и снова



\[

F[f * g]=\sqrt{2 \pi} F[f] F[g] .

\]



С. обобщенных функций широко используется при решении краевых задач для уравнений с частными производными. Так, интеграл Пуассона, написанный в виде



\[

U(x, t)=\mu(x) * \frac{1}{2 \sqrt{\pi t}} e^{-x^{2} / 4 t},

\]



дает решение уравнения тешлошроводности для бесконечного стержня, когда начальная температура $\mu(x)$ может быть не только обычной, но и обобщенной функцией.



Понятие С. как обычных, так и обобщенных функций естественным образом переносится на случай функций многих независимых переменных: надо в предыдущем считать $x, y$ не действительными числами, а векторами из $\mathbb{R} n$.



Лит.: [1] В л а д и м и р о в В. С., Уравнения математической физики, 4 изд., М., 1981; [2] Г е л ь ф а н д И. М., ІІІ и ло о в Г. Е., Обобщенные Функции и действия над ними, М., 1958; пзер. с и т ч м а р ш Е., Введение в теорию интегралов Фурье,

п. И. Соболев.





\end{document}